% Vietnamese translation for OI Public Library
% Source-File: IOI中国国家候选队论文/2003/何林/何林.ppt
\documentclass[11pt]{article}
\usepackage{oipl-vn}

\title{Bản trình chiếu: Một cách giải cho lớp bài toán cân bi}
\author{He Lin}
\date{2003}

\begin{document}
\maketitle

\origtitle{Slide companion for ball-weighing problems}
\sourcepath{IOI China candidate team papers / 2003 / He Lin / He Lin.ppt}

\section{Phạm vi bản dịch}

Đây là bản trình chiếu đi kèm bài viết về bài toán cân bi bằng cân hai đĩa.
Các slide khôi phục được chứa nhiều công thức dạng \(\log_3 n\), các trường
hợp \(n=3k+1\), \(n=3k+2\), và các nhóm \(A_1,A_2,A_3\). Nội dung dưới đây
tóm lược mạch trình bày chính bằng tiếng Việt.

\section{Cận dưới}

Mỗi lần cân có ba kết quả, vì vậy một cây quyết định độ sâu \(h\) phân biệt
tối đa \(3^h\) trạng thái. Do đó số lần cân cần thiết thỏa
\[
  h \ge \lceil \log_3 n\rceil
\]
hoặc biến thể tương ứng khi số trạng thái cần phân biệt là \(2n\).

\section{Xây dựng phương án}

Tư tưởng chính là chia tập bi thành ba phần gần bằng nhau. Với
\(n=3k+1\) hoặc \(n=3k+2\), các nhóm được phân bố thành kích thước
\(k,k+1\) sao cho sau một lần cân, số khả năng còn lại ở các nhánh không quá
lớn. Các công thức trong slide kiểm tra rằng độ sâu còn lại đạt cận
\(\lceil\log_3(2n)\rceil\) trong các trường hợp biên.

\section{Liên hệ với bài viết}

Bản bài viết đã dịch trình bày đầy đủ hơn bằng cây quyết định và quy nạp.
Bản trình chiếu là phiên bản rút gọn, tập trung vào các công thức chia nhóm
và các ví dụ \(N=3\), \(N=9\).

\end{document}
